\section{OperatorIdentities}

\subsection{Commutator Bracket}

\begin{definition}[opComm]
The \textbf{commutator} (Lie bracket) of two $\R$-linear operators $f, g : J \to J$ is:
\[ \opComm{f}{g} \defas f \circ g - g \circ f \]
\end{definition}

\begin{notation}
We write $\llbracket f, g \rrbracket$ for $\opComm{f}{g}$.
\end{notation}

\begin{lemma}[opComm\_apply]
For all $\R$-linear maps $f, g : J \to J$ and $x \in J$:
\[ \llbracket f, g \rrbracket(x) = f(g(x)) - g(f(x)) \]
\end{lemma}

\begin{lemma}[opComm\_skew]
For all $\R$-linear maps $f, g : J \to J$:
\[ \llbracket f, g \rrbracket = -\llbracket g, f \rrbracket \]
\end{lemma}

\begin{lemma}[opComm\_self]
For all $\R$-linear maps $f : J \to J$:
\[ \llbracket f, f \rrbracket = 0 \]
\end{lemma}

\subsection{Linearized Jordan Identity}

\begin{theorem}[linearized\_jordan\_operator]
For all $a, b, c \in J$:
\[ 2 \cdot \bigl( \llbracket \Lop{a}, \Lop{b \jmul c} \rrbracket + \llbracket \Lop{b}, \Lop{c \jmul a} \rrbracket + \llbracket \Lop{c}, \Lop{a \jmul b} \rrbracket \bigr) = 0 \]
\end{theorem}

\begin{theorem}[linearized\_on\_jsq]
For all $a, b, c \in J$:
\begin{align*}
2 \cdot \bigl( &a \jmul ((b \jmul c) \jmul \jsq{a}) - (b \jmul c) \jmul (a \jmul \jsq{a}) \\
+ &b \jmul ((c \jmul a) \jmul \jsq{a}) - (c \jmul a) \jmul (b \jmul \jsq{a}) \\
+ &c \jmul ((a \jmul b) \jmul \jsq{a}) - (a \jmul b) \jmul (c \jmul \jsq{a}) \bigr) = 0
\end{align*}
\end{theorem}

\begin{theorem}[linearized\_core]
For all $a, b, c \in J$:
\begin{align*}
&a \jmul ((b \jmul c) \jmul \jsq{a}) - (b \jmul c) \jmul (a \jmul \jsq{a}) \\
+ &b \jmul ((c \jmul a) \jmul \jsq{a}) - (c \jmul a) \jmul (b \jmul \jsq{a}) \\
+ &c \jmul ((a \jmul b) \jmul \jsq{a}) - (a \jmul b) \jmul (c \jmul \jsq{a}) = 0
\end{align*}
\end{theorem}

\begin{theorem}[linearized\_rearranged]
For all $a, b, c \in J$:
\begin{align*}
&a \jmul ((b \jmul c) \jmul \jsq{a}) + b \jmul ((c \jmul a) \jmul \jsq{a}) + c \jmul ((a \jmul b) \jmul \jsq{a}) \\
&= (b \jmul c) \jmul (a \jmul \jsq{a}) + (c \jmul a) \jmul (b \jmul \jsq{a}) + (a \jmul b) \jmul (c \jmul \jsq{a})
\end{align*}
\end{theorem}

\subsection{Operator Commutator with Squared Element}

\begin{theorem}[operator\_commutator\_jsq]
For all $a, b \in J$:
\[ \llbracket \Lop{\jsq{a}}, \Lop{b} \rrbracket = 2 \cdot \llbracket \Lop{a}, \Lop{a \jmul b} \rrbracket \]
\end{theorem}

\begin{theorem}[operator\_commutator\_jsq\_apply]
For all $a, b, x \in J$:
\[ \jsq{a} \jmul (b \jmul x) - b \jmul (\jsq{a} \jmul x) = 2 \cdot \bigl( a \jmul ((a \jmul b) \jmul x) - (a \jmul b) \jmul (a \jmul x) \bigr) \]
\end{theorem}

\subsection{Idempotent Operator Identities}

\begin{theorem}[opComm\_idempotent\_eigenspace]
Let $e \in J$ be idempotent. Let $a \in J$ with $e \jmul a = \lambda \cdot a$ for some $\lambda \in \R$. Then for all $x \in J$:
\[ \llbracket \Lop{e}, \Lop{a} \rrbracket(x) = e \jmul (a \jmul x) - a \jmul (e \jmul x) \]
\end{theorem}

\begin{theorem}[L\_e\_L\_a\_L\_e]
Let $e \in J$ be idempotent. For all $a \in J$:
\[ \Lop{e} \circ \Lop{a} \circ \Lop{e} = \Lop{e \jmul (a \jmul e)} + \frac{1}{2} \cdot \llbracket \Lop{e}, \llbracket \Lop{e}, \Lop{a} \rrbracket \rrbracket \]
\emph{(Currently unproven in the formalization.)}
\end{theorem}

\begin{theorem}[opComm\_double\_idempotent]
Let $e \in J$ be idempotent. For all $a \in J$:
\[ \llbracket \Lop{e}, \llbracket \Lop{e}, \Lop{a} \rrbracket \rrbracket = 2 \cdot (\Lop{e} \circ \Lop{a} \circ \Lop{e}) - 2 \cdot \Lop{e \jmul (a \jmul e)} \]
\emph{(Currently unproven in the formalization.)}
\end{theorem}

\subsection{Useful Commutator Properties}

\begin{theorem}[opComm\_on\_P1]
Let $e \in J$ be idempotent. Let $a, b \in J$ with $e \jmul a = a$ and $e \jmul b = b$ (i.e., $a, b \in \Peirce{e}{1}$). Then:
\[ \llbracket \Lop{e}, \Lop{a} \rrbracket(b) = e \jmul (a \jmul b) - a \jmul b \]
\end{theorem}

\begin{theorem}[opComm\_on\_P0]
Let $e \in J$ be idempotent. Let $a \in J$ with $e \jmul a = 0$ (i.e., $a \in \Peirce{e}{0}$). Then for all $x \in J$:
\[ \llbracket \Lop{e}, \Lop{a} \rrbracket(x) = e \jmul (a \jmul x) - a \jmul (e \jmul x) \]
\end{theorem}
